\documentclass{article}
\usepackage{graphicx} % Required for inserting images

% added stuff
\usepackage{amssymb}
\usepackage{amsmath}

\newcommand{\R}{\mathbb{R}}
\newcommand{\N}{\mathbb{N}}
\newcommand{\Z}{\mathbb{Z}}
\newcommand{\Q}{\mathbb{Q}}
\newcommand{\C}{\mathbb{C}}

\newcommand{\power}{\mathcal{P}}
\newcommand{\T}{\mathcal{T}}
\newcommand{\B}{\mathcal{B}}

\newcommand{\eps}{\varepsilon}
\newcommand{\nullset}{\varnothing}
\newcommand{\lb}{\left(}
\newcommand{\rb}{\right)}
% end added stuff

\title{Fundamentals 21}
\author{Siddharth K}
\date{April 2025}

\begin{document}

\maketitle

\subsection*{Problem 21.1}
Let $\T = \{\nullset, X\}$. We show that it satisfies the three axioms:
\begin{enumerate}
    \item $\nullset, X \in \T$.
    \item $\nullset \cup X = X \in \T$.
    \item $\nullset \cap X = \nullset \in \T$.
\end{enumerate}
Now, let $\T = \power(X)$.
\begin{enumerate}
    \item $\nullset$ and $X$ are subsets of $X$ so $\nullset, X \in \power(X)$.
    \item Assume $\{X_\alpha: \alpha \in A\} \subseteq \T$, then for any $\alpha$ and $x \in X_\alpha$, $x \in X$. So $\bigcup_{\alpha \in A} X_\alpha \in \power(X)$.
    \item Let $Y, Z \in \power(X)$, so that $Y$ and $Z$ are subsets of $X$. We have that $Y \cap Z$ is a subset of $X$, so $Y\cap Z \in \power(X)$.
\end{enumerate}

\subsection*{Problem 21.3}
Since the evenly spaced integer topology is defined as the set of unions of all $\{a\Z + b\}$, it suffices to show that $\B = \{a\Z + b : a, b \in \Z, a \neq 0\}$ is a topological basis.

\begin{enumerate}
    \item For any $x \in \Z$, $x \in \{1\Z + x\} \in \B$.
    \item If $x \in a\Z + b$ and $c\Z + d$, then $x \in \text{lcm}(a, c)\Z + x \subseteq (a\Z + b) \cap (c\Z + d)$.
\end{enumerate}

\subsection*{Problem 21.4}
First, we show that $\T$ is a topological basis:

\begin{enumerate}
    \item For any $x \in X$, since $X \in \T$, $x \in X \in \T$.
    \item If $x \in B_1 \cap B_2$, since $B_3 = B_1 \cap B_2 \in \T$, $x \in B_3 \subseteq B_1\cap B_2$.
\end{enumerate}

Since $\T$ is closed under arbitrary unions, all unions of elements in $\T$ are also in $\T$, so $\T$ can be generated by itself.

\subsection*{Problem 21.5}
First assume that $\B$ is a basis for $\T$. Let $U \in \T$ and $x \in U$. We know that $U$ is a union of elements in $\B$, say $\{B_\alpha : \alpha \in A\}$. Since $x \in U$, we have that $x \in B_\alpha$ for some $\alpha \in A$, and $B_\alpha \subseteq U$ because $U$ is a union of $B_\alpha$ and other sets. \\

Now, assume that for each $U \in \T$ and $x \in U$ there's $B \in \B$ such that $x \in B \subseteq U$. We need to show that all elements of $\T$ are unions of elements in $\B$. For contradiction, assume this is false. Then there's $U \in \T$ where $U \neq \bigcup_{\alpha \in A} X_\alpha$ for any $\{X_\alpha : \alpha \in A\} \subseteq \B$. Then there's some $x \in U$ where $x \not \in \B$ for any $B \in \B$, which is a contradiction.

\subsection*{Problem 21.6}
\subsubsection*{Part (1)}
\begin{enumerate}
    \item Let $x \in \R$ and $B = (x - 1, x + 1)$. Then $B \in \B$ and $x \in B$.
    \item Let $B_3 = B_1 \cap B_2$. Since $B_3$ is open, $B_3 \in \B$.
\end{enumerate}

\subsubsection*{Part (2)}
We first show that every open set in $\R$ is a union of open intervals. Let $U \in \R$ open. Then for all $x \in X$, there exists $\eps_x$ such that $B_\eps(x) \subseteq U$. Therefore, $\bigcup_{x \in X} B_{\eps_x}(x) \subseteq U$. Also, since $x$ is in some $B_{\eps_x}(x)$ for all $x \in U$, $U \subseteq \bigcup_{x \in X} B_{\eps_x}(x)$. So we have that $U = \bigcup_{x \in X} B_{\eps_x}(x)$. \\

Now, since all open sets can be written as arbitrary unions of elements in $\B$, the topology generated by $\B$ is just all sets that are open by the Euclidean metric. \\

\subsubsection*{Part (3)}
Every open set is a union of open intervals $(a, b)$, so it suffices to show that all open intervals are arbitrary unions of intervals $(a, b)$ where $a, b \in \Q$. We can find rational sequences $(a_n), (b_n)$ such that $a = \lim_{n\to\infty} a_n$ and $b = \lim_{n\to\infty} b_n$, where $a_n > a$ and $b_n < b$ for all $n \in \N$. So we have that $(a, b) = \bigcup_{n \in \N} (a_n, b_n)$.

\subsection*{Problem 21.8}
Let $\{X_\alpha : \alpha \in A\} \subseteq X$ be a set of neighborhoods of $x$. So there is some $U \in \T$ where $x \in U \subseteq X_\alpha$ for some $\alpha \in A$. Then $U \subseteq \bigcup_{\alpha \in A} X_\alpha$, so $\bigcup_{\alpha \in A} X_\alpha$ is a neighborhood containing $X$. \\

Let $Y, Z \subseteq X$ be neighborhoods of $x$. Then there are $U_1 \subseteq Y$, $U_2 \subseteq Z$ where $U_1, U_2 \in \T$. So $U_1 \cap U_2 \subseteq Y \cap Z$, and since $U_1\cap U_2$ is open and $x \in U_1 \cap U_2$, $Y\cap Z$ is a neighborhood of $x$.

\subsection*{Problem 21.9}
Let $T = \bigcap_{\alpha \in A} \T_\alpha$.
\begin{enumerate}
    \item $\nullset, X \in \T_\alpha$ for all $\alpha \in A$, s $\nullset, X \in \T$.
    \item Let $B$ be any set. If $\{X_\beta: \beta \in B\} \subseteq \T$ then $\{X_\beta: \beta \in B\} \in \T_\alpha$ for all $\alpha \in A$. Since arbitrary unions of sets in $\T_\alpha$ are also in $\T_\alpha$ for all $\alpha \in A$, $\bigcup_{\beta \in B} X_\beta \in \T_\alpha$ for all $\alpha \in A$, so $\bigcup_{\beta \in B} X_\beta \in \T$.
    \item Let $Y, Z \in \T$, so that $Y, Z \in \T_\alpha$ for all $\alpha \in A$. We have that $Y \cap Z \in T_\alpha$ for all $\alpha \in A$, so $Y \cap Z \in \T$.
\end{enumerate}

\subsection*{Problem 21.10}
We can find a counterexample. Let $\T$ be the topology induced by the Euclidean metric on $\R$ and $\T'$ be the topology induced by the discrete metric on $\R$. Open balls in $\T'$ are one element sets. So, for any $x \in \R$, the set ${x} \in T'$ is contained in $(x - 1, x + 1) \in \T$. Also, for any $U \in \T$, $U$ is contained in the set $\bigcup_{x \in U} \{x\} \in T'$. \\

We have that the set $\{0, 1\} \in \T'$, but it is not open under the Euclidean metric, so $\{0, 1\} \not\in \T$.

\end{document}
